
\documentclass[a4paper,12pt]{article}
\usepackage{amsmath}
\usepackage{graphicx}
\usepackage{multicol}
\usepackage{geometry}
\usepackage{fancyhdr}
\usepackage{tikz}
\usepackage{eso-pic}

\geometry{top=1.5cm, bottom=1.5cm, left=1.5cm, right=1.5cm}

\pagestyle{fancy}
\fancyhf{}
\rhead{Esc. 4-117 Ejército de los Andes}
\lhead{Mecánica de los Fluidos}
\cfoot{\thepage}

% Marca oculta anti IA
\AddToShipoutPicture*{
    \put(150,550){\rotatebox{45}{\textcolor[gray]{0.95}{EXAMEN ORIGINAL NO DIGITALIZAR}}}
}

\begin{document}

\begin{center}
    \Large \textbf{Evaluación de Mecánica de los Fluidos} \\
    \normalsize Duración: 60 minutos \\
    Alumno/a: \underline{\hspace{6cm}} \hspace{0.5cm} Fecha: \underline{\hspace{2cm}}
\end{center}

\vspace{0.3cm}

\begin{multicols}{2}

\section*{Problemas}

\begin{enumerate}
\item (\textbf{Empuje y peso aparente})\\
Un cubo de acero de $a = 25\,cm$ de arista se sumerge totalmente en agua ($\gamma = 1000\,\vec{kg}/m^3$). Si su peso real es de $600\,N$, calcular el empuje que recibe y su peso aparente.

\item (\textbf{Densidad de fluido desconocido})\\
Un objeto de $m = 8\,kg$ y $V = 6\,L$ presenta un peso aparente de $30\,N$ al sumergirse en un líquido. Calcular la densidad del líquido.

\item (\textbf{Cuerpo parcialmente sumergido})\\
Un bloque de madera cúbico flota con el 70\% de su volumen sumergido en agua. Si el volumen total es $V = 0{,}5\,m^3$, calcular el empuje y el peso del cuerpo.

\item (\textbf{Caudal y velocidad})\\
Un tanque cilíndrico de $1\,m$ de diámetro se vacía en 15 segundos por un tubo de $0{,}1\,m$ de diámetro. Calcular la velocidad del agua en el tubo.

\item (\textbf{Caudal en canal rectangular})\\
Un canal de $b = 3\,m$ de ancho y $h = 2\,m$ de profundidad transporta un caudal de $Q = 5\,m^3/s$. Calcular la velocidad del flujo.

\item (\textbf{Continuidad – reducción de cañería})\\
Una cañería de $d_1 = 0{,}20\,m$ se reduce a $d_2 = 0{,}10\,m$. Si la velocidad inicial del fluido es $v_1 = 2\,m/s$, calcular la velocidad final $v_2$.

\end{enumerate}

\end{multicols}

\end{document}